\section{AHSME 1964}
        \begin{enumerate}
        	\item (\textit{AHSME 1964}) What is the value of $[\log_{10}(5\log_{10}100)]^2$?


 $\textbf{(A)}\ \log_{10}50 \qquad \textbf{(B)}\ 25\qquad \textbf{(C)}\ 10 \qquad \textbf{(D)}\ 2\qquad \textbf{(E)}\ 1$


	\item (\textit{AHSME 1964}) The graph of $x^2-4y^2=0$ is:


 $\textbf{(A)}\ \text{a parabola} \qquad \textbf{(B)}\ \text{an ellipse} \qquad \textbf{(C)}\ \text{a pair of straight lines}\qquad \\ \textbf{(D)}\ \text{a point}\qquad \textbf{(E)}\ \text{None of these}$


	\item (\textit{AHSME 1964}) When a positive integer $x$ is divided by a positive integer $y$, the quotient is $u$ and the remainder is $v$, where $u$ and $v$ are integers. 
What is the remainder when $x+2uy$ is divided by $y$?


 $\textbf{(A)}\ 0 \qquad \textbf{(B)}\ 2u \qquad \textbf{(C)}\ 3u \qquad \textbf{(D)}\ v \qquad \textbf{(E)}\ 2v$


	\item (\textit{AHSME 1964}) The expression


 
 \[\frac{P+Q}{P-Q}-\frac{P-Q}{P+Q}\]


 where $P=x+y$ and $Q=x-y$, is equivalent to:


 $\textbf{(A)}\ \frac{x^2-y^2}{xy}\qquad \textbf{(B)}\ \frac{x^2-y^2}{2xy}\qquad \textbf{(C)}\ 1 \qquad \textbf{(D)}\ \frac{x^2+y^2}{xy}\qquad \textbf{(E)}\ \frac{x^2+y^2}{2xy}$


	\item (\textit{AHSME 1964}) If $y$ varies directly as $x$, and if $y=8$ when $x=4$, the value of $y$ when $x=-8$ is:


 $\textbf{(A)}\ -16 \qquad \textbf{(B)}\ -4 \qquad \textbf{(C)}\ -2 \qquad \textbf{(D)}\ 4k, k= \pm1, \pm2, \dots \qquad \\ \textbf{(E)}\ 16k, k=\pm1,\pm2,\dots$


	\item (\textit{AHSME 1964}) If $x, 2x+2, 3x+3, \dots$ are in geometric progression, the fourth term is:


 $\textbf{(A)}\ -27 \qquad \textbf{(B)}\ -13\frac{1}{2} \qquad \textbf{(C)}\ 12\qquad \textbf{(D)}\ 13\frac{1}{2}\qquad \textbf{(E)}\ 27$


	\item (\textit{AHSME 1964}) Let n be the number of real values of $p$ for which the roots of 
$x^2-px+p=0$
are equal. Then n equals:


 $\textbf{(A)}\ 0 \qquad \textbf{(B)}\ 1 \qquad \textbf{(C)}\ 2 \qquad \textbf{(D)}\ \text{a finite number greater than 2}\qquad \textbf{(E)}\ \infty$


	\item (\textit{AHSME 1964}) The smaller root of the equation 
$\left(x-\frac{3}{4}\right)\left(x-\frac{3}{4}\right)+\left(x-\frac{3}{4}\right)\left(x-\frac{1}{2}\right) =0$ is:


 $\textbf{(A)}\ -\frac{3}{4}\qquad \textbf{(B)}\ \frac{1}{2}\qquad \textbf{(C)}\ \frac{5}{8}\qquad \textbf{(D)}\ \frac{3}{4}\qquad \textbf{(E)}\ 1$


	\item (\textit{AHSME 1964}) A jobber buys an article at $\$24$ less $12\frac{1}{2}\%$. He then wishes to sell the article at a gain of $33\frac{1}{3}\%$ of his cost 
after allowing a $20\%$ discount on his marked price. At what price, in dollars, should the article be marked?


 $\textbf{(A)}\ 25.20 \qquad \textbf{(B)}\ 30.00 \qquad \textbf{(C)}\ 33.60 \qquad \textbf{(D)}\ 40.00 \qquad \textbf{(E)}\ \text{none of these}$


	\item (\textit{AHSME 1964}) Given a square side of length $s$. On a diagonal as base a triangle with three unequal sides is constructed so that its area
equals that of the square. The length of the altitude drawn to the base is:


 $\textbf{(A)}\ s\sqrt{2} \qquad \textbf{(B)}\ s/\sqrt{2} \qquad \textbf{(C)}\ 2s \qquad \textbf{(D)}\ 2\sqrt{s} \qquad \textbf{(E)}\ 2/\sqrt{s}$


	\item (\textit{AHSME 1964}) Given $2^x=8^{y+1}$ and $9^y=3^{x-9}$, find the value of $x+y$
$\textbf{(A)}\ 18 \qquad \textbf{(B)}\ 21 \qquad \textbf{(C)}\ 24 \qquad \textbf{(D)}\ 27 \qquad \textbf{(E)}\ 30$


	\item (\textit{AHSME 1964}) Which of the following is the negation of the statement: For all $x$ of a certain set, $x^2>0$?


 $\textbf{(A)}\ \text{For all x}, x^2 < 0\qquad \textbf{(B)}\ \text{For all x}, x^2 \le 0\qquad \textbf{(C)}\ \text{For no x}, x^2>0\qquad \\ \textbf{(D)}\ \text{For some x}, x^2>0\qquad \textbf{(E)}\ \text{For some x}, x^2 \le 0$


	\item (\textit{AHSME 1964}) A circle is inscribed in a triangle with side lengths $8, 13$, and $17$. Let the segments of the side of length $8$, 
made by a point of tangency, be $r$ and $s$, with $r<s$. What is the ratio $r:s$? 


 $\textbf{(A)}\ 1:3 \qquad \textbf{(B)}\ 2:5 \qquad \textbf{(C)}\ 1:2 \qquad \textbf{(D)}\ 2:3 \qquad \textbf{(E)}\ 3:4$


	\item (\textit{AHSME 1964}) A farmer bought $749$ sheep. He sold $700$ of them for the price paid for the $749$ sheep. 
The remaining $49$ sheep were sold at the same price per head as the other $700$. 
Based on the cost, the percent gain on the entire transaction is:


 $\textbf{(A)}\ 6.5 \qquad \textbf{(B)}\ 6.75 \qquad \textbf{(C)}\ 7 \qquad \textbf{(D)}\ 7.5 \qquad \textbf{(E)}\ 8$


	\item (\textit{AHSME 1964}) A line through the point $(-a,0)$ cuts from the second quadrant a triangular region with area $T$. The equation of the line is:


 $\textbf{(A)}\ 2Tx+a^2y+2aT=0 \qquad \textbf{(B)}\ 2Tx-a^2y+2aT=0 \qquad \textbf{(C)}\ 2Tx+a^2y-2aT=0 \qquad \\ \textbf{(D)}\ 2Tx-a^2y-2aT=0 \qquad \textbf{(E)}\ \text{none of these}$


	\item (\textit{AHSME 1964}) Let $f(x)=x^2+3x+2$ and let $S$ be the set of integers $\{0, 1, 2, \dots , 25 \}$. 
The number of members $s$ of $S$ such that $f(s)$ has remainder zero when divided by $6$ is:


 $\textbf{(A)}\ 25\qquad \textbf{(B)}\ 22\qquad \textbf{(C)}\ 21\qquad \textbf{(D)}\ 18 \qquad \textbf{(E)}\ 17$


	\item (\textit{AHSME 1964}) Given the distinct points $P(x_1, y_1), Q(x_2, y_2)$ and $R(x_1+x_2, y_1+y_2)$. 
Line segments are drawn connecting these points to each other and to the origin $O$. 
Of the three possibilities: (1) parallelogram (2) straight line (3) trapezoid, figure $OPRQ$, 
depending upon the location of the points $P, Q$, and $R$, can be:


 $\textbf{(A)}\ \text{(1) only}\qquad \textbf{(B)}\ \text{(2) only}\qquad \textbf{(C)}\ \text{(3) only}\qquad \textbf{(D)}\ \text{(1) or (2) only}\qquad \textbf{(E)}\ \text{all three}$


	\item (\textit{AHSME 1964}) Let $n$ be the number of pairs of values of $b$ and $c$ such that $3x+by+c=0$ and $cx-2y+12=0$ have the same graph. Then $n$ is:


 $\textbf{(A)}\ 0\qquad \textbf{(B)}\ 1\qquad \textbf{(C)}\ 2\qquad \textbf{(D)}\ \text{finite but more than 2}\qquad \textbf{(E)}\ \infty$


	\item (\textit{AHSME 1964}) If $2x-3y-z=0$ and $x+3y-14z=0, z \neq 0$, the numerical value of $\frac{x^2+3xy}{y^2+z^2}$ is:


 $\textbf{(A)}\ 7\qquad \textbf{(B)}\ 2\qquad \textbf{(C)}\ 0\qquad \textbf{(D)}\ -20/17\qquad \textbf{(E)}\ -2$


	\item (\textit{AHSME 1964}) The sum of the numerical coefficients of all the terms in the expansion of $(x-2y)^{18}$ is:


 $\textbf{(A)}\ 0\qquad \textbf{(B)}\ 1\qquad \textbf{(C)}\ 19\qquad \textbf{(D)}\ -1\qquad \textbf{(E)}\ -19$


	\item (\textit{AHSME 1964}) If $\log_{b^2}x+\log_{x^2}b=1, b>0, b \neq 1, x \neq 1$, then $x$ equals:


 $\textbf{(A)}\ 1/b^2 \qquad \textbf{(B)}\ 1/b \qquad \textbf{(C)}\ b^2 \qquad \textbf{(D)}\ b \qquad \textbf{(E)}\ \sqrt{b}$


	\item (\textit{AHSME 1964}) Given parallelogram $ABCD$ with $E$ the midpoint of diagonal $BD$. Point $E$ is connected to a point $F$ in $DA$ so that 
$DF=\frac{1}{3}DA$. What is the ratio of the area of $\triangle DFE$ to the area of quadrilateral $ABEF$?


 $\textbf{(A)}\ 1:2 \qquad \textbf{(B)}\ 1:3 \qquad \textbf{(C)}\ 1:5 \qquad \textbf{(D)}\ 1:6 \qquad \textbf{(E)}\ 1:7$


	\item (\textit{AHSME 1964}) Two numbers are such that their difference, their sum, and their product are to one another as $1:7:24$. The product of the two numbers is:


 $\textbf{(A)}\ 6\qquad \textbf{(B)}\ 12\qquad \textbf{(C)}\ 24\qquad \textbf{(D)}\ 48\qquad \textbf{(E)}\ 96$


	\item (\textit{AHSME 1964}) Let $y=(x-a)^2+(x-b)^2, a, b$ constants. For what value of $x$ is $y$ a minimum?


 $\textbf{(A)}\ \frac{a+b}{2} \qquad \textbf{(B)}\ a+b \qquad \textbf{(C)}\ \sqrt{ab} \qquad \textbf{(D)}\ \sqrt{\frac{a^2+b^2}{2}}\qquad \textbf{(E)}\ \frac{a+b}{2ab}$


	\item (\textit{AHSME 1964}) The set of values of $m$ for which $x^2+3xy+x+my-m$ has two factors, with integer coefficients, which are linear in $x$ and $y$, is precisely:


 $\textbf{(A)}\ 0, 12, -12\qquad \textbf{(B)}\ 0, 12\qquad \textbf{(C)}\ 12, -12\qquad \textbf{(D)}\ 12\qquad \textbf{(E)}\ 0$


	\item (\textit{AHSME 1964}) In a ten-mile race $\textit{First}$ beats $\textit{Second}$ by $2$ miles and 
$\textit{First}$ beats $\textit{Third}$ by $4$ miles. 
If the runners maintain constant speeds throughout the race, 
by how many miles does $\textit{Second}$ beat $\textit{Third}$?


 $\textbf{(A)}\ 2\qquad \textbf{(B)}\ 2\frac{1}{4}\qquad \textbf{(C)}\ 2\frac{1}{2}\qquad \textbf{(D)}\ 2\frac{3}{4}\qquad \textbf{(E)}\ 3$


	\item (\textit{AHSME 1964}) If $x$ is a real number and $|x-4|+|x-3|<a$ where $a>0$, then:


 $\textbf{(A)}\ 0<a<.01\qquad \textbf{(B)}\ .01<a<1 \qquad \textbf{(C)}\ 0<a<1\qquad \\ \textbf{(D)}\ 0<a \le 1\qquad \textbf{(E)}\ a>1$


	\item (\textit{AHSME 1964}) The sum of $n$ terms of an arithmetic progression is $153$, and the common difference is $2$. 
If the first term is an integer, and $n>1$, then the number of possible values for $n$ is:


 $\textbf{(A)}\ 2\qquad \textbf{(B)}\ 3\qquad \textbf{(C)}\ 4\qquad \textbf{(D)}\ 5\qquad \textbf{(E)}\ 6$


	\item (\textit{AHSME 1964}) In this figure $\angle RFS = \angle FDR, FD = 4$ inches, $DR = 6$ inches, $FR = 5$ inches, $FS = 7\frac{1}{2}$ inches. 
The length of $RS$, in inches, is:


 
                    \begin{figure}[H]
                        \centering
                        \includegraphics[width=0.5\textwidth]{Problem/AHSME_I_1964_29_Problem_Figure_1.png}
                    \end{figure}
                    
$\textbf{(A)}\ \text{undetermined} \qquad \textbf{(B)}\ 4\qquad \textbf{(C)}\ 5\frac{1}{2} \qquad \textbf{(D)}\ 6\qquad \textbf{(E)}\ 6\frac{1}{4}$


	\item (\textit{AHSME 1964}) If $(7+4\sqrt{3})x^2+(2+\sqrt{3})x-2=0$, the larger root minus the smaller root is:


 $\textbf{(A)}\ -2+3\sqrt{3}\qquad \textbf{(B)}\ 2-\sqrt{3}\qquad \textbf{(C)}\ 6+3\sqrt{3}\qquad \textbf{(D)}\ 6-3\sqrt{3}\qquad \textbf{(E)}\ 3\sqrt{3}+2$


	\item (\textit{AHSME 1964}) Let $f(n)=\dfrac{5+3\sqrt{5}}{10}\left(\dfrac{1+\sqrt{5}}{2}\right)^n+\dfrac{5-3\sqrt{5}}{10}\left(\dfrac{1-\sqrt{5}}{2}\right)^n$. 
Then $f(n+1)-f(n-1)$, expressed in terms of $f(n)$, equals:


 $\textbf{(A)}\ \dfrac{1}{2}f(n) \qquad \textbf{(B)}\ f(n)\qquad \textbf{(C)}\ 2f(n)+1 \qquad \textbf{(D)}\ f^2(n)\qquad \textbf{(E)}\ \frac{1}{2}(f^2(n)-1)$


	\item (\textit{AHSME 1964}) If $\dfrac{a+b}{b+c}=\dfrac{c+d}{d+a}$, then:


 $\textbf{(A)}\ a\text{ must equal }c \qquad \textbf{(B)}\ a+b+c+d\text{ must equal zero }\qquad \textbf{(C)}\ \text{either }a=c\text{ or }a+b+c+d=0,\text{ or both} \qquad \textbf{(D)}\ a+b+c+d\neq 0\text{ if }a=c \qquad \\ \textbf{(E)}\ a(b+c+d)=c(a+b+d)$


	\item (\textit{AHSME 1964}) $P$ is a point interior to rectangle $ABCD$ and such that $PA=3$ inches, $PD=4$ inches, and $PC=5$ inches. Then $PB$, in inches, equals:


 $\textbf{(A) }2\sqrt{3}\qquad \textbf{(B) }3\sqrt{2}\qquad \textbf{(C) }3\sqrt{3}\qquad \textbf{(D) }4\sqrt{2}\qquad \textbf{(E) }2$

                    \begin{figure}[H]
                        \centering
                        \includegraphics[width=0.5\textwidth]{Problem/AHSME_I_1964_33_Problem_Figure_1.png}
                    \end{figure}
                    


	\item (\textit{AHSME 1964}) If $n$ is a multiple of $4$, the sum $s=1+2i+3i^2+ \ldots +(n+1)i^{n}$, where $i=\sqrt{-1}$, equals:


 $\textbf{(A)}\ 1+i\qquad \textbf{(B)}\ \frac{1}{2}(n+2) \qquad \textbf{(C)}\ \frac{1}{2}(n+2-ni) \qquad \\ \textbf{(D)}\ \frac{1}{2}[(n+1)(1-i)+2]\qquad \textbf{(E)}\ \frac{1}{8}(n^2+8-4ni)$


	\item (\textit{AHSME 1964}) The sides of a triangle are of lengths $13, 14,$ and $15$. The altitudes of the triangle meet at point $H$. 
If $AD$ is the altitude to the side length $14$, what is the ratio $HD:HA$?     


 

	\item (\textit{AHSME 1964}) In this figure the radius of the circle is equal to the altitude of the equilateral triangle $ABC$. 
The circle is made to roll along the side $AB$, remaining tangent to it at a variable point $T$ 
and intersecting lines $AC$ and $BC$ in variable points $M$ and $N$, respectively. 
Let $n$ be the number of degrees in arc $MTN$. Then $n$, for all permissible positions of the circle:


 $\textbf{(A) }\text{varies from }30^{\circ}\text{ to }90^{\circ}\quad \textbf{(B) }\text{varies from }30^{\circ}\text{ to }60^{\circ} \\ \textbf{(C) }\text{varies from }60^{\circ}\text{ to }90^{\circ} \quad \textbf{(D) }\text{remains constant at }30^{\circ}\quad \textbf{(E) }\text{remains constant at }60^{\circ}$

                    \begin{figure}[H]
                        \centering
                        \includegraphics[width=0.5\textwidth]{Problem/AHSME_I_1964_36_Problem_Figure_1.png}
                    \end{figure}
                    


	\item (\textit{AHSME 1964}) Given two positive number $a, b$ such that $a<b$ , let $A.M.$ be their arithmetic mean and let $G.M.$ be their positive geometric mean. 
Then $A.M.$ minus $G.M.$ is always less than:


 $\textbf{(A) }\dfrac{(b+a)^2}{ab}\qquad \textbf{(B) }\dfrac{(b+a)^2}{8b}\qquad \textbf{(C) }\dfrac{(b-a)^2}{ab} \textbf{(D) }\dfrac{(b-a)^2}{8a}\qquad  \textbf{(E) }\dfrac{(b-a)^2}{8b}$


	\item (\textit{AHSME 1964}) The sides $PQ$ and $PR$ of $\triangle PQR$ are respectively of lengths $4$ inches, and $7$ inches. The median $PM$ is $3\frac{1}{2}$ inches. 
Then $QR$, in inches, is:


 $\textbf{(A) }6\qquad \textbf{(B) }7\qquad \textbf{(C) }8\qquad \textbf{(D) }9\qquad  \textbf{(E) }10$


	\item (\textit{AHSME 1964}) The magnitudes of the sides of $\triangle ABC$ are $a, b$, and $c$, as shown, with $c\le b\le a$. 
Through interior point $P$ and the vertices $A, B, C$, lines are drawn meeting the opposite sides in $A', B', C'$, respectively. 
Let $s=AA'+BB'+CC'$. Then, for all positions of point $P$, $s$ is less than:


 $\textbf{(A) }2a+b\qquad \textbf{(B) }2a+c\qquad \textbf{(C) }2b+c\qquad \textbf{(D) }a+2b\qquad  \textbf{(E) } a+b+c$

                    \begin{figure}[H]
                        \centering
                        \includegraphics[width=0.5\textwidth]{Problem/AHSME_I_1964_39_Problem_Figure_1.png}
                    \end{figure}
                    


	\item (\textit{AHSME 1964}) A watch loses $2\frac{1}{2}$ minutes per day. It is set right at $1$ P.M. on March $15$. 
Let $n$ be the positive correction, in minutes, to be added to the time shown by the watch at a given time. 
When the watch shows $9$ A.M. on March $21$, $n$ equals:


 $\textbf{(A) }14\frac{14}{23}\qquad \textbf{(B) }14\frac{1}{14}\qquad \textbf{(C) }13\frac{101}{115}\qquad \textbf{(D) }13\frac{83}{115}\qquad \textbf{(E) }13\frac{13}{23}$


\end{enumerate}